\scanpage{137}
\mspage{13}
and $\dim\bar R=n$; $x_1,\ldots,x_n$ is also a minimal basis of $\bar M$. Then
for any non-negative integer $q$, $y\in M^q$ and $y\notin M^{q+1}$ implies
that $y=f(x_1,\ldots,x_n)$ where $f(X_1,\ldots,X_n)$ is a form of
degree $q$ with coefficients in $R$ but not all in $M$; suppose if
possible that $y\in\bar M^{q+1}$, then $y=g(x_1,\ldots,x_n)$ where
$g(X_1,\ldots,X_n)$ is a form of degree $q+1$ with coefficients in
$\bar R$, so that $y=h(x_1,\ldots,x_n)$ where $h(X_1,\ldots,X_n)$ is a form
of degree $q$ with coefficients in $\bar M$, let $n(X_1,\ldots,X_n)=$
$f(X_1,\ldots,X_n)-h(X_1,\ldots,X_n)$, since $f$ has coefficients in $R$
but not all in $M$ and since $\bar M\cap R=M$, these coefficients are
in $\bar R$ but not all in $\bar M$, and since the coefficients of $h$ are
all in $\bar M$ we conclude that $n(X_1,\ldots,X_n)$ is a form of degree
$q$ with coefficients in $\bar R$ but not all in $\bar M$ and hence
$n(x_1,\ldots,x_n)\notin\bar M^{q+1}$ so that $n(x_1,\ldots,x_n)\ne0$ which is a
contradiction since $n(x_1,\ldots,x_n)=y-y=0$; consequently,
$y\in M^q$ and $y\notin M^{q+1}$ implies that $y\notin\bar M^{q+1}$. Therefore, for
any positive integer $q$, $y\in R$ and $y\notin M^q$ implies that $y\notin\bar M^q$;
and hence $y\in R$ and $y\in\bar M^q$ implies that $y\in M^q$. Therefore
$R\cap\bar M^q\subset M^q$ and hence $R\cap\bar M^q=M^q$.

Next, let $A$ be any ideal in $R$. Let $y\in(A\bar R)\cap R$ and let
$q$ be any positive integer; then $y\in A\bar R$ implies that $y=\Sigma a_i z_i$
with $a_i$ in $A$ and $z_i$ in $\bar R$ so that $z_i=u_i+t_i$ with $u_i$
in $R$ and $t_i$ in $\bar M^q$, hence $y=\Sigma a_i u_i+\Sigma a_i t_i$ with $\Sigma a_i u_i$ in
$A$ and $\Sigma a_i t_i$ in $\bar M^q$; since $\Sigma a_i t_i=y-\Sigma a_i u_i$ with $y$ in $R$
and $\Sigma a_i u_i$ in $R$ we have $\Sigma a_i t_i\in R\cap\bar M^q=M^q$ so that
$y=\Sigma a_i u_i+\Sigma a_i t_i\in A+M^q$. Thus $(A\bar R)\cap R\subset\deleted{\unsure{\bigcap_{q=1}^{\infty}}=A}\inserted{\bigcap_{q=1}^{\infty}(A+M^q)=A}$ and hence
\figure{\begin{tikzpicture}\draw[->] (1,0) to[out=180,in=-90] (0,1);\end{tikzpicture}}
\mspage{136}
